% Generated by scripts/assemble.py; corrections belong in source/corrections.json.
\examhead{186}{长安大学 2026 年数学分析试卷}{175}
% p176-b24; source PDF page 176
\par\Needspace{5\baselineskip}
\blockanchor{p176-b24}
\majorquestion{一}计算题. 前 4 题每题 12 分, 第 5 题 14 分, 共 62 分.\par
% p176-b25; source PDF page 176
\blockanchor{p176-b25}
\question{1}计算极限  $\lim_{x\to0}\frac{\cos(\sqrt{2}x)-e^{-x^{2}}+\frac{x^{4}}{3}}{x^{6}}$ .\par
% p176-b26; source PDF page 176
\blockanchor{p176-b26}
\question{2}设函数 $y = f(x)$ 由方程 $y^3 + xy^2 + x^2y + 6 = 0$ 确定，求 $f(x)$ 的极值.\par
% p176-b27; source PDF page 176
\blockanchor{p176-b27}
\question{3}计算积分 $J = \int_0^1 \ln \frac{1 + x}{1 + x^2} \, \mathrm{d}x$.\par
% p176-b28; source PDF page 176
\blockanchor{p176-b28}
\question{4}计算重积分 $\iint_{D} y \sin(x + y) \, \mathrm{d}x \, \mathrm{d}y$, 其中 $D = \{(x, y) \mid 0 \leq x + y \leq \frac{\pi}{2}, 0 \leq x - y \leq \frac{\pi}{2}\}$.\par
% p177-b00; source PDF page 177
\par\Needspace{6\baselineskip}
\blockanchor{p177-b00}
\question{5}已知有向曲线 $L$ 为球面 $x^{2} + y^{2} + z^{2} = 2x$ 与平面 $2x - z - 1 = 0$ 的交线, 从 $z$ 轴正向往 $z$ 轴负向看为逆时针方向, 计算曲线积分\par
\nopagebreak[4]
% p177-b01; source PDF page 177
\blockanchor{p177-b01}
\displayformula{\oint_ {L} \left(6 x y z - y z ^ {2}\right) \mathrm{d} x + 2 x ^ {2} z \mathrm{d} y + x y z \mathrm{d} z.}
% p177-b02; source PDF page 177
\par\Needspace{5\baselineskip}
\blockanchor{p177-b02}
\majorquestion{二}证明题. 每题 12 分, 共 48 分.\par
% p177-b03; source PDF page 177
\par\Needspace{6\baselineskip}
\blockanchor{p177-b03}
\question{6}设 $S$ 为有界数集, 证明: 若 $\sup S = a \notin S$, 则存在严格单调递增的数列 $\{x_n\} \subset S$, 使得\par
\nopagebreak[4]
% p177-b04; source PDF page 177
\blockanchor{p177-b04}
\displayformula{\lim _ {n \to \infty} x _ {n} = a.}
% p177-b05; source PDF page 177
\blockanchor{p177-b05}
\question{7}设函数 $f(x), g(x)$ 在 $[a, b]$ 上连续，在 $(a, b)$ 内具有二阶导数且存在相等的最大值，$f(a) = g(a)$，$f(b) = g(b)$，证明：存在 $\xi \in (a, b)$，使得 $f''(\xi) = g''(\xi)$.\par
% p177-b06; source PDF page 177
\par\Needspace{6\baselineskip}
\blockanchor{p177-b06}
\question{8}证明: 若 $f(x)$ 在 $[a, b]$ 上可积, $F(x)$ 在 $[a, b]$ 上连续, 且除有限个点外有 $F'(x) = f(x)$, 则有\par
\nopagebreak[4]
% p177-b07; source PDF page 177
\blockanchor{p177-b07}
\displayformula{\int_ {a} ^ {b} f (x) \mathrm{d} x = F (b) - F (a).}
% p177-b08; source PDF page 177
\blockanchor{p177-b08}
\question{9}证明: 若 $f(x)$ 在 $[a, +\infty)$ 上可导, 且 $\int_{a}^{+\infty} f(x) \mathrm{d}x$ 与 $\int_{a}^{+\infty} f'(x) \mathrm{d}x$ 都收敛, 则 $\lim_{x \to +\infty} f(x) = 0$.\par
% p177-b09; source PDF page 177
\par\Needspace{5\baselineskip}
\blockanchor{p177-b09}
\majorquestion{三}综合题. 每题 20 分, 共 40 分.\par
% p177-b10; source PDF page 177
\blockanchor{p177-b10}
\question{10}解答如下问题: 设 $f(x) = \sum_{n=2}^{\infty} \frac{x^n}{n^2 \ln n}$. 证明:\par
% p177-b11; source PDF page 177
\blockanchor{p177-b11}
(1) 证明: $f(x)$ 在 $[-1, 1]$ 上连续.\par
% p177-b12; source PDF page 177
\blockanchor{p177-b12}
(2) 证明: $f(x)$ 在 $x = -1$ 处右可导.\par
% p177-b13; source PDF page 177
\blockanchor{p177-b13}
(3) 证明: $\lim_{x \to 1^{-}} f'(x) = +\infty$.\par
% p177-b14; source PDF page 177
\blockanchor{p177-b14}
(4) 证明: $f(x)$ 在 $x = 1$ 处左不可导.\par
% p177-b15; source PDF page 177
\blockanchor{p177-b15}
\question{11}解答如下问题: 设 $I(k, a) = \int_{0}^{+\infty} e^{-kx} \frac{\sin(ax)}{x} \mathrm{d}x$.\par
% p177-b16; source PDF page 177
\blockanchor{p177-b16}
(1) 证明: 对固定的 $k \in [0, +\infty)$, 积分 $I(k, a)$ 关于 $a$ 在 $|a| \geq \delta (\delta > 0)$ 上一致收敛, 对固定的 $a \neq 0$, 积分 $I(k, a)$ 关于 $k$ 在 $[0, +\infty)$ 上一致收敛.\par
% p177-b17; source PDF page 177
\blockanchor{p177-b17}
(2) 计算积分 $I = \int_{0}^{+\infty} \frac{\sin x}{x} \, \mathrm{d}x$ 的值.\par
